\section{Scope, finite expectations, and abstract state functionals}
\label{sec:ch19-scope}

A correlation function is an expectation of an \emph{ordered product} of
observables.  In elementary quantum mechanics this usually means inserting
operators between a normalized vector and its adjoint, or taking their trace
against a density matrix.  For the bosonization problem developed here, that
familiar operation must be distinguished from the evaluation of words in an
\emph{abstract} canonical-commutation-relation (CCR) algebra.  The distinction
is essential: the finite fermionic Hilbert space has genuine normalized states,
but it cannot carry nonzero scalar CCR as identities of all its operators.
Conversely, the uncompressed oscillator algebra admits an exactly positive,
Gaussian vacuum functional, but the construction of this functional does not
identify it with any finite-band fermionic ground state.

The prerequisites are the finite CAR representation and occupation-state
inner product of Chapter~\ref{ch:car_fock}, the finite energy budgets of
Chapter~\ref{ch:vacuum_budget}, and the polynomial oscillator construction of
Chapter~\ref{ch:boson_fock}.  We use the dressed-mode transformation from
Chapter~18, with real parameters obeying $c^2-s^2=1$.  No infinite-volume
limit, Hilbert-space completion of the polynomial oscillator carrier, or
identification of finite fermionic and abstract bosonic expectations is assumed.
The aim of this section is to construct and distinguish the two kinds of
state evaluation.  We explicitly distinguish the finite CAR carrier (which models actual fermionic lattice states), the ambient polynomial oscillator carrier (unbounded degree, before completion), the protected budget (a subspace with valid commutation properties), the algebraic state (the abstract expectation functional $\omega_P$), and the coefficientwise formal-series evaluation used in later models. Ordered mode contractions, their spatial Fourier sums,
and quantitative comparisons between the models are treated subsequently.
For the usual physical context of bosonized correlators, we refer the reader to standard analytical treatments such as
\cite{giamarchi2004,vondelft1998}. The exact finite hypotheses used below
are those of the present construction rather than consequences of continuum
bosonization prescriptions.

\subsection{Expectations in a finite quantum system}

Let $\mathcal H_{\mathrm{fin}}$ be a nonzero, finite-dimensional complex
Hilbert space.  We use the convention that $\langle\phi,\psi\rangle$ is
conjugate-linear in its first argument and linear in its second.  Its
observable algebra is $\mathcal B(\mathcal H_{\mathrm{fin}})$, the algebra of
all linear operators on this space, with the Hilbert adjoint $X^\dagger$.
A \emph{normalized vector state} is specified by
$\psi\in\mathcal H_{\mathrm{fin}}$ with $\langle\psi,\psi\rangle=1$:
\begin{equation}
  \langle X\rangle_{\psi,\mathrm{fin}}
   :=\langle\psi,X\psi\rangle,
  \qquad X\in\mathcal B(\mathcal H_{\mathrm{fin}}).
  \label{eq:ch19-vector-expectation}
\end{equation}
More generally, a \emph{density matrix} is an operator $\Gamma$ satisfying
$\Gamma=\Gamma^\dagger$, $\langle v,\Gamma v\rangle\geq0$ for all $v$,
and $\operatorname{Tr}\Gamma=1$.  Its expectation is
\begin{equation}
  \langle X\rangle_{\Gamma,\mathrm{fin}}
     :=\operatorname{Tr}(\Gamma X).
  \label{eq:ch19-density-expectation}
\end{equation}
The pure state in \eqref{eq:ch19-vector-expectation} is recovered by taking
$\Gamma=|\psi\rangle\langle\psi|$.  Both functionals are complex-linear,
normalized, and positive: for example,
\begin{align}
  \langle I\rangle_{\Gamma,\mathrm{fin}}&=1,\\
  \langle X^\dagger X\rangle_{\psi,\mathrm{fin}}
    &=\|X\psi\|^2\geq0,\\
  \langle X^\dagger X\rangle_{\Gamma,\mathrm{fin}}
    &=\operatorname{Tr}\bigl((X\Gamma^{1/2})^\dagger
                                  X\Gamma^{1/2}\bigr)\geq0.
  \label{eq:ch19-finite-positivity}
\end{align}
In particular, a state is a linear functional, \emph{not} generally a
multiplicative map: the number $\langle XY\rangle$ cannot be replaced by
$\langle X\rangle\langle Y\rangle$ without an additional hypothesis.
Moreover, the order of $X$ and $Y$ matters if they do not commute.

A self-adjoint finite Hamiltonian $H_{\mathrm{fin}}$ has a smallest
eigenvalue $E_0$ and at least one normalized eigenvector $\psi_0$ with
$H_{\mathrm{fin}}\psi_0=E_0\psi_0$.  Its ground expectation is therefore
well-defined by \eqref{eq:ch19-vector-expectation}; any positive trace-one
matrix supported on the ground eigenspace also defines a ground-state
expectation.  This elementary existence result does \emph{not} say that the
same vector is killed by operators obtained by compressing a Bogoliubov
transformation.  The latter is an independent, potentially inconsistent set
of equations.

In the finite CAR setting of the earlier chapters,
$\mathcal H_{\mathrm{fin}}$ may be the occupation Fock space of a finite
momentum band, or an energy-budget subspace with compressed operators.
Whenever we write a finite expectation we specify its vector or density
matrix and the operators acting on the corresponding space.  In particular,
these brackets will never denote the abstract oscillator functional
$\omega_P$ defined below.

\subsection{Why an exact scalar CCR cannot hold on a finite space}

\begin{proposition}[Trace obstruction to finite-dimensional scalar CCR]
\label{prop:ch19-trace-obstruction}
Let $\mathcal H$ be a nonzero finite-dimensional complex vector space,
$A,C\in\operatorname{End}_{\mathbb C}(\mathcal H)$, and
$\lambda\in\mathbb C$.  If
\begin{equation}
  [A,C]:=AC-CA=\lambda I_{\mathcal H},
  \label{eq:ch19-claimed-finite-ccr}
\end{equation}
then $\lambda=0$.  In particular, $[A,C]=mI$ is impossible for a
positive integer $m$.
\end{proposition}
\begin{proof}
Writing $d=\dim_{\mathbb C}\mathcal H>0$ and using cyclicity of the
ordinary matrix trace, we obtain
\[
 \operatorname{Tr}[A,C]
   =\operatorname{Tr}(AC)-\operatorname{Tr}(CA)=0.
\]
On the other hand, \eqref{eq:ch19-claimed-finite-ccr} gives
$\operatorname{Tr}[A,C]=\lambda\operatorname{Tr}I=\lambda d$.
Since the coefficient field has characteristic zero and $d>0$,
$\lambda=0$.
\end{proof}

It is important to state exactly what the obstruction excludes.  Finite
fermionic density operators can have precisely computed commutators on
selected states, sectors, or energy budgets, and these calculations may
reproduce the expected central term on those states.  What cannot hold is
a nonzero \emph{scalar operator identity on the entire finite space}.
Compression need not preserve a CCR: if $P$ is a finite-rank orthogonal
projection in an ambient representation, $A_P=PAP|_{P\mathcal H}$, and
$C_P=PCP|_{P\mathcal H}$, then on a common domain on which all products
make sense,
\begin{equation}
 [A_P,C_P]
 =\left(P[A,C]P-PA(I-P)CP+PC(I-P)AP\right)|_{P\mathcal H}.
 \label{eq:ch19-compression-remainder}
\end{equation}
The last two terms are boundary contributions: they retain the excursions
outside the projected space that are lost when operators are multiplied
after compression.  Even if $[A,C]=mI$ in the ambient algebra, the
compressed commutator cannot equal $mI_{P\mathcal H}$ for $m>0$ and
$0<\operatorname{rank}P<\infty$.  For the finite CAR densities considered
earlier, the analogous residuals must be estimated under their actual
band, margin, and excursion hypotheses; they are not removed by simply
calling the finite currents ``bosons.''

\subsection{An uncompressed CCR algebra and its polynomial vacuum}

\paragraph{Why an uncompressed description is necessary.}
In the harmonic oscillator of undergraduate quantum mechanics, the raising
operator can be applied repeatedly:
$|0\rangle,|1\rangle,|2\rangle,\ldots$.  Even one oscillator has no
largest occupation number.  The algebraic model below retains precisely this
feature without first introducing an infinite series or a Hilbert-space
completion.  We shall work with a \emph{finite set of oscillator modes}, but
allow polynomials of \emph{arbitrarily large finite degree} in their creation
operators.  ``Uncompressed'' therefore refers to the absence of a degree or
energy cutoff, \emph{not} to infinitely many momentum labels, infinite
lattice size, or an implicit thermodynamic limit.

This distinction matters because a finite-dimensional truncation can retain
many low-occupation matrix elements of an oscillator and still fail its
operator commutation relations.  The trace obstruction just proved applies
to the truncated matrices, whereas the uncompressed differential operators
below act on an infinite-dimensional vector space.  Our goal is to establish
an exact algebraic setting for the CCR before assigning expectations to
operator products.

\paragraph{Three objects that should not be confused.}
We shall use three elementary notions from algebra.  An \emph{algebra} over
$\mathbb C$ is a complex vector space with an associative, bilinear product;
we require a multiplicative identity $I$.  For example, matrices form an
algebra under matrix multiplication, while polynomials form an algebra under
ordinary multiplication.  An element of an algebra need not yet be a
concrete linear operator.  A \emph{representation} supplies that extra step:
it is a map taking abstract algebra elements to operators while respecting
sums, scalar multiplication, products, and the identity.  In particular,
representation of a product is composition, so the rightmost operator acts
first.

The \emph{carrier} of a mathematical structure means the underlying set or
vector space on which its operations are defined.  Here there are two
relevant carriers: the vector space of formal operator expressions forming
the abstract algebra, and the polynomial vector space on which these
expressions will act.  We will call the latter the \emph{representation
carrier} $\mathcal P_M$.  Neither word denotes a spatial lattice, a
single-particle momentum band, or a finite-particle Fock space.  The term
\emph{ambient} likewise means ``the larger space or algebra before a
specified restriction.''  In this construction the ambient word algebra
precedes imposition of the CCR, while the uncompressed polynomial space is
ambient to its finite-degree subspaces.  The two uses will always be
identified by their symbols.

\paragraph{Free words and the relations we wish to impose.}
Fix an integer $M\geq1$ and the finite set of mode labels
\[
  \mathcal I_M=\{(\nu,m):\nu\in\{R,L\},\;1\leq m\leq M\}.
\]
For each $j=(\nu,m)$ introduce two \emph{formal symbols}: $A_j$, intended
for annihilation, and $C_j$, intended for creation.  Before any relations
have been imposed, an expression such as $A_j C_k A_\ell$ is simply an
ordered \emph{word}; it is not yet a matrix or differential operator.
Consider all finite complex linear combinations of such words, including
the empty word $I$.  Their multiplication is concatenation followed by
linear extension.  This is the \emph{free unital algebra}, which we denote
by $\mathfrak F_M$.  In $\mathfrak F_M$ the words $A_jC_k$ and $C_kA_j$
are distinct, just as two products of noncommuting quantum observables
must be kept in their stated order.

The word algebra also has an \emph{involution}, denoted by $*$: it exchanges
$A_j$ with $C_j$, conjugates complex coefficients, and reverses every
word.  More explicitly,
\begin{equation}
 (A_j)^*=C_j,\qquad (C_j)^*=A_j,\qquad
 (\lambda u+\mu v)^*=\overline\lambda u^*+
       \overline\mu v^*,\qquad (uv)^*=v^*u^*.
 \label{eq:ch19-free-star}
\end{equation}
This is the algebraic counterpart of the Hilbert-space adjoint.  It becomes
an actual adjoint relation only after we construct a suitable inner product
and representation.

The canonical commutation relations should hold \emph{as algebraic
equalities}, independently of any particular vector.  For clarity define
the desired relation expressions
\begin{align}
 r^{AC}_{jk}&:=A_jC_k-C_kA_j-\delta_{jk}m_j I,\nonumber\\
 r^{AA}_{jk}&:=A_jA_k-A_kA_j,\qquad
 r^{CC}_{jk}:=C_jC_k-C_kC_j,
 \label{eq:ch19-relation-expressions}
\end{align}
where $j=(\nu,m_j)$ and $\delta_{jk}=1$ exactly when both the branch and
the positive mode index agree.  The expressions $r^{AC}_{jk}$,
$r^{AA}_{jk}$, and $r^{CC}_{jk}$ are generally \emph{not zero in the
free algebra}.  They are the relations that we now impose.

\paragraph{What an ideal and a quotient do.}
An \emph{ideal} $\mathfrak J$ of an algebra is a linear subspace that is
closed under multiplication by arbitrary algebra elements on the left
\emph{and} on the right: if $r\in\mathfrak J$ and $u,v\in\mathfrak F_M$,
then $urv\in\mathfrak J$.  This two-sided requirement is crucial for
operator algebra.  If $A_jC_k-C_kA_j=\delta_{jk}m_j I$, the same equality
must remain valid when inserted inside a longer word, for example after
multiplying by $C_\ell$ on the left and $A_p$ on the right.

Let $\mathfrak J_{\mathrm{CCR}}$ be the smallest two-sided ideal
containing all the expressions in
\eqref{eq:ch19-relation-expressions} and their involutions.  Thus its
elements are finite linear combinations of expressions $urv$ built from
these relations (including their stars).  The inclusion of involutions
makes the ideal \emph{$*$-stable}: $r\in\mathfrak J_{\mathrm{CCR}}$
implies $r^*\in\mathfrak J_{\mathrm{CCR}}$.  This ensures that taking
adjoints is compatible with the equations being imposed.

The \emph{quotient algebra}
\begin{equation}
 \mathfrak A_{\mathrm{CCR}}(M)
   :=\mathfrak F_M/\mathfrak J_{\mathrm{CCR}}
 \label{eq:ch19-ccr-quotient}
\end{equation}
identifies two word expressions $u$ and $v$ when
$u-v\in\mathfrak J_{\mathrm{CCR}}$.  We are not setting every word to
zero; only expressions implied by the selected relations become zero.
For instance, in the quotient $A_jC_j=C_jA_j+m_j I$, but neither
$A_jC_j$ nor $C_jA_j$ is required to vanish.  This is a systematic way
to turn formal commutation rules into exact identities, while preserving
the order of all remaining products.

The resulting unital $*$-algebra is characterized by
\begin{align}
 [A_{\nu m},C_{\eta n}]
   &=\delta_{\nu\eta}\delta_{mn}\,mI,\nonumber\\
 [A_{\nu m},A_{\eta n}]&=0,
 \qquad [C_{\nu m},C_{\eta n}]=0.
 \label{eq:ch19-abstract-ccr}
\end{align}
These equalities hold in $\mathfrak A_{\mathrm{CCR}}(M)$ for \emph{all}
finite words, not merely after taking a vacuum expectation.
One must still check that the quotient has not collapsed to the zero
algebra (equivalently, that $I$ has not been identified with zero), and
that it admits a \emph{positive} evaluation.  Both will follow from a
concrete polynomial representation.

\paragraph{From abstract words to an explicit representation carrier.}
The variables $X_j$ introduced next are \emph{commuting polynomial
variables}, unlike the noncommuting $A_j,C_j$ in $\mathfrak F_M$.
The purpose of the representation is not to equate these two algebras:
it sends the abstract generators to \emph{operators on polynomials}.
Creation becomes multiplication by a variable, whereas annihilation
becomes differentiation.  Even though the polynomial variables commute,
multiplication and differentiation generally do not, which is precisely
how the CCR arise.  A representation satisfying the relations annihilates
every element of $\mathfrak J_{\mathrm{CCR}}$ and therefore gives a
well-defined action of the quotient on its carrier.

Let $\mathcal P_M=\mathbb C[X_j:j\in\mathcal I_M]$, the space of
polynomials in commuting indeterminates.  For a multi-index
$\alpha=(\alpha_j)_{j\in\mathcal I_M}$ with $\alpha_j\in\mathbb N$,
write $X^\alpha=\prod_jX_j^{\alpha_j}$.  Give this vector space the
Hermitian form
\begin{equation}
 \langle X^\alpha,X^\beta\rangle_{\mathcal P}
 =\delta_{\alpha\beta}
     \prod_{(\nu,m)\in\mathcal I_M}\alpha_{\nu m}!\,
                                           m^{\alpha_{\nu m}}.
 \label{eq:ch19-polynomial-gram}
\end{equation}
Every weight in the product is strictly positive.  Thus if
$p=\sum_\alpha p_\alpha X^\alpha$ is nonzero, then
\begin{equation}
 \|p\|_{\mathcal P}^2
 =\sum_\alpha |p_\alpha|^2
    \prod_{(\nu,m)}\alpha_{\nu m}!\,m^{\alpha_{\nu m}}>0.
 \label{eq:ch19-polynomial-norm}
\end{equation}
Only finitely many terms occur in each polynomial; no completion or
convergence of infinite sums is involved.  Define linear operators on
$\mathcal P_M$ by
\begin{equation}
  \pi(C_{\nu m})p=X_{\nu m}p,
  \qquad
  \pi(A_{\nu m})p=m\frac{\partial p}{\partial X_{\nu m}}.
  \label{eq:ch19-polynomial-representation}
\end{equation}
They are adjoint on this polynomial pre-Hilbert space.  To verify this
without appealing to an infinite-dimensional operator theorem, fix
$j=(\nu,m)$ and monomials $X^\alpha,X^\beta$.  The only potentially
nonzero matrix element of $\pi(A_j)$ from $X^\alpha$ to $X^\beta$ occurs
when $\alpha=\beta+e_j$.  In that case,
\[
 \langle X^\beta,\pi(A_j)X^\alpha\rangle_{\mathcal P}
   =m\alpha_j\,\|X^\beta\|_{\mathcal P}^2
   =\|X^\alpha\|_{\mathcal P}^2
   =\langle\pi(C_j)X^\beta,X^\alpha\rangle_{\mathcal P}.
\]
Thus $\pi(A_j)^\dagger=\pi(C_j)$ in the algebraic sense that these
pairings agree on every pair of polynomials.  The elementary identity
$\partial_{X_j}(X_kp)-X_k\partial_{X_j}p=\delta_{jk}p$ gives
\[
 [\pi(A_{\nu m}),\pi(C_{\eta n})]
     =\delta_{\nu\eta}\delta_{mn}\,mI_{\mathcal P_M},
\]
and the two homogeneous commutators vanish.  Because the CCR hold on every polynomial, the representation of the free
word algebra sends each relation expression, including its left and
right multiples, to zero.  It therefore factors through the quotient,
with the unit represented by the nonzero identity on $\mathcal P_M$.
In particular, the quotient is not the zero algebra.  Faithfulness of
the representation is not needed to establish this conclusion.

\paragraph{One-mode calculation: where the missing top level goes.}
For a single fixed label $j=(\nu,m)$, suppress the subscripts and consider
$A=m\,d/dX$ and $C$ acting by multiplication by $X$.  On every monomial
$X^n$, $n\geq0$,
\begin{align}
 AC\,X^n&=m(n+1)X^n, & CA\,X^n&=mnX^n,\\
 [A,C]X^n&=mX^n.&&
 \label{eq:ch19-one-mode-ccr}
\end{align}
Thus the commutator is exactly $mI$ on the entire polynomial carrier,
including the vacuum $1=X^0$.  Repeated creation yields
$C^n1=X^n$ and
\begin{equation}
 \|C^n1\|_{\mathcal P}^2=n!\,m^n,
 \qquad n=0,1,2,\ldots .
 \label{eq:ch19-one-mode-norms}
\end{equation}
No polynomial has infinite degree, but for \emph{every} $n$ there is a
nonzero vector of degree $n$.  This is the elementary reason that the
representation space has infinite dimension even when only one mode is
considered.

For comparison, keep only $1,X,\ldots,X^N$ and \emph{compress} both
operators to that $(N+1)$-dimensional subspace.  The derivative still
lowers degree, but the compressed creator kills $X^N$ because $X^{N+1}$
has been discarded.  Consequently the compressed commutator has action
\begin{equation}
 [A_N,C_N]X^n=
 \begin{cases}
 mX^n,&0\leq n<N,\\
 -mN X^N,&n=N.
 \end{cases}
 \label{eq:ch19-top-degree-defect}
\end{equation}
The failure of $mI$ is concentrated at the top degree, and the trace is
$mN-mN=0$, exactly as required by
Proposition~\ref{prop:ch19-trace-obstruction}.  This example exhibits
both the usefulness and the limitation of an energy cutoff: low-degree
vectors may obey the expected commutator while the global finite matrix
identity is impossible.  The trace obstruction does not apply to
$\mathcal P_M$ itself, since it is infinite-dimensional; moreover, the
algebraic multiplication and differentiation maps are not being asserted
to be bounded operators on a Hilbert completion.

\begin{definition}[Polynomial vacuum functional]
\label{def:ch19-polynomial-vacuum}
For $a\in\mathfrak A_{\mathrm{CCR}}(M)$ set
\begin{equation}
   \omega_0(a):=\langle1,\pi(a)1\rangle_{\mathcal P}.
   \label{eq:ch19-omega-zero}
\end{equation}
The vector $1$ is the constant polynomial, of norm one.
\end{definition}

\begin{proposition}[Normalization, positivity, and vacuum orientation]
\label{prop:ch19-polynomial-state}
The linear functional $\omega_0$ satisfies
\begin{align}
 \omega_0(I)&=1,\qquad
 \omega_0(a^*a)=\|\pi(a)1\|_{\mathcal P}^2\geq0,
 \label{eq:ch19-omega0-positive}\\
 \omega_0(aA_{\nu m})&=0,\qquad
 \omega_0(C_{\nu m}a)=0
 \quad\text{for every }a\in\mathfrak A_{\mathrm{CCR}}(M).
 \label{eq:ch19-omega0-annihilation}
\end{align}
In particular, the state is not a multiplicative algebra homomorphism.
\end{proposition}
\begin{proof}
Linearity and normalization follow from the polynomial representation
and $\|1\|_{\mathcal P}=1$.  Algebraic adjointness gives
$\langle1,\pi(a^*a)1\rangle_{\mathcal P}
 =\langle\pi(a)1,\pi(a)1\rangle_{\mathcal P}$,
which proves positivity.  For the first vanishing equation,
$\pi(A_{\nu m})1=0$; the corresponding operator occurs on the
\emph{right} of the product, so it acts first.  For the second,
adjointness gives
\[
 \langle1,\pi(C_{\nu m})\pi(a)1\rangle_{\mathcal P}
 =\langle\pi(A_{\nu m})1,\pi(a)1\rangle_{\mathcal P}=0.
\]
Finally, for any $m\geq1$,
\[
 \omega_0(A_{\nu m})=\omega_0(C_{\nu m})=0,
 \quad\omega_0(A_{\nu m}C_{\nu m})=m,
 \quad\omega_0(C_{\nu m}A_{\nu m})=0.
\]
The nonzero two-point expectation is inconsistent with
multiplicativity of a map having zero one-point values.
\end{proof}

The polynomial realization also explains the term \emph{Gaussian} or
\emph{quasi-free}.  Expectations of any finite word reduce by successively
commuting annihilators past creators, until a rightmost annihilator or a
leftmost creator kills the vacuum expectation.  The recursion expresses
even moments as sums of pairings of generators, while odd moments vanish.
Equivalently, all higher truncated moments beyond degree two vanish in this
algebraic oscillator state.  This is a statement about finite polynomial
words; no analytic generating-function convergence has been assumed.

\subsection{Dressing the vacuum by a Bogoliubov automorphism}

Let $P$ stand for real transformation parameters $(c,s)$ satisfying
\begin{equation}
 c,s\in\mathbb R,\qquad c^2-s^2=1.
 \label{eq:ch19-hyperbolic}
\end{equation}
The larger Luttinger-parameter record used in the preceding chapter can
provide such a pair, but no separate value of a Luttinger parameter is
needed for the present construction.  Write $\bar R=L$, $\bar L=R$.
Define on the generators
\begin{align}
 \beta_P(C_{\nu m})&=cC_{\nu m}+sA_{\bar\nu m},
 \label{eq:ch19-bogoliubov-creator}\\
 \beta_P(A_{\nu m})&=cA_{\nu m}+sC_{\bar\nu m}.
 \label{eq:ch19-bogoliubov-annihilator}
\end{align}
The parameters are real, so these formulas preserve the involution.
Their mixed commutator is multiplied by $c^2-s^2=1$; their
same-type commutators still vanish (the cross-orientation terms cancel).
Thus $\beta_P$ preserves the defining relations
\eqref{eq:ch19-abstract-ccr} and extends to a star-algebra
automorphism of $\mathfrak A_{\mathrm{CCR}}(M)$.  Direct substitution
shows that
\begin{equation}
 \beta_P^{-1}=\beta_{(c,-s)}.
 \label{eq:ch19-inverse-beta}
\end{equation}
We call $\widetilde A_{\nu m}:=\beta_P(A_{\nu m})$ and
$\widetilde C_{\nu m}:=\beta_P(C_{\nu m})$ the \emph{dressed}
generators.  Notice that a Bogoliubov transformation of abstract
CCR generators is not, by itself, a unitary transformation on the
finite CAR Hilbert space.

\begin{definition}[Dressed abstract state]
\label{def:ch19-dressed-functional}
The dressed oscillator state is the linear functional
\begin{equation}
  \boxed{\displaystyle
  \omega_P:=\omega_0\circ\beta_P^{-1}}
  \quad\text{on }\mathfrak A_{\mathrm{CCR}}(M).
  \label{eq:ch19-dressed-state}
\end{equation}
\end{definition}

\begin{proposition}[Positive dressed vacuum]
\label{prop:ch19-dressed-positive}
Under \eqref{eq:ch19-hyperbolic}, $\omega_P$ is normalized and
positive, and for every $a\in\mathfrak A_{\mathrm{CCR}}(M)$,
\begin{equation}
 \omega_P(a\widetilde A_{\nu m})=0,
 \qquad
 \omega_P(\widetilde C_{\nu m}a)=0.
 \label{eq:ch19-dressed-annihilation}
\end{equation}
\end{proposition}
\begin{proof}
Because $\beta_P^{-1}$ is a unital star-algebra automorphism,
\[
 \omega_P(I)=\omega_0(I)=1,\qquad
 \omega_P(a^*a)
 =\omega_0\!\left((\beta_P^{-1}a)^*\beta_P^{-1}a\right)\geq0.
\]
For the first annihilation identity, multiplicativity of the
\emph{automorphism}, not of the state, gives
\begin{align*}
 \omega_P(a\widetilde A_{\nu m})
 &=\omega_0\!\left(\beta_P^{-1}(a)
                      \beta_P^{-1}(\beta_P(A_{\nu m}))\right)\\
 &=\omega_0(\beta_P^{-1}(a)A_{\nu m})=0.
\end{align*}
Similarly,
$\omega_P(\widetilde C_{\nu m}a)
 =\omega_0(C_{\nu m}\beta_P^{-1}(a))=0$.
The direction of the pullback is therefore necessary:
using $\omega_0\circ\beta_P$ instead would not, in general, yield the
vacuum annihilated by the dressed lowering generators just defined.
\end{proof}

The familiar undressed case is $c=1$, $s=0$, for which
$\beta_P=\mathrm{id}$ and $\omega_P=\omega_0$.  A nonzero value of $s$
changes the covariance of the bare modes, not the abstract CCR
itself.  The dressed functional is again quasi-free because substituting
linear combinations of $A$ and $C$ into any finite word preserves the
finite-word Gaussian reduction described above.  Nevertheless, it is not
an ordinary matrix trace on $\mathcal H_{\mathrm{fin}}$, and it must not
be denoted by $\langle\cdot\rangle_{\psi,\mathrm{fin}}$.

\paragraph{Relation to the accepted Lean construction.}
The implementation in \texttt{A08Oscillator}\allowbreak\texttt{States.lean}
represents the \emph{free} word algebra by
\texttt{OscAlgebra M}, the polynomials by \texttt{OscPoly M},
and the creators and annihilators by \texttt{oscCreation} and
\texttt{osc}\allowbreak\texttt{Annihilation}.  The functionals \texttt{omega0} and
\texttt{omegaP} are complex-linear maps on that free-word carrier,
with $\texttt{omegaP}=\texttt{omega0}\circ
\texttt{bogoliubovMap}(c,-s)$.  Theorems such as
\texttt{omega0\_}\allowbreak\texttt{gramSum}, \texttt{omegaP\_}\allowbreak\texttt{gramSum},
\texttt{omegaP\_ccr\_}\allowbreak\texttt{descends} and its homogeneous-relation variants
establish positivity of finite word combinations and invariance under
insertion of the CCR relations; \texttt{omegaP\_comp\_}\allowbreak\texttt{bogoliubovMap}
and the dressed annihilation theorems fix the orientation of the
pullback.  The quotient $\mathfrak A_{\mathrm{CCR}}(M)$ above is a
mathematical presentation of the relation algebra, \emph{not} a
separately instantiated Lean quotient type in A08.  In that
implementation, the ambient \texttt{FreeAlgebra} remains the carrier of
\texttt{omega0} and \texttt{omegaP}.  The explicit theorems showing
that their values are unchanged when any defining relation is inserted
between arbitrary words justify descending these functionals to the
corresponding relation quotient.  This distinction between a
\emph{presented quotient} and an \emph{evaluation proved to respect its
relations} is a genuine interface distinction, not merely a matter of
notation.
Likewise, the implementation's \texttt{B\_complete} status records
acceptance of that development; it is not promotion to the frozen Core.

\subsection{Worked finite example: compression can destroy the dressed vacuum}

We now return to an ordinary three-dimensional Hilbert space.  The
example makes the separation of \eqref{eq:ch19-density-expectation}
and \eqref{eq:ch19-dressed-state} concrete.  Retain the two modes
$(R,1)$ and $(L,1)$, but project the polynomial carrier to total
degree at most one:
\[
 \mathcal H_{\leq1}
  :=\operatorname{span}_{\mathbb C}\{|0\rangle,|R\rangle,|L\rangle\}
  =\operatorname{span}_{\mathbb C}\{1,X_R,X_L\}.
\]
For $m=1$ the Gram weights in \eqref{eq:ch19-polynomial-gram}
make the displayed ordered basis orthonormal.  Let $D_R,D_L$ be the
compressed derivatives and $\widehat X_R,\widehat X_L$ the compressed
multiplication operators.  Consider the compressed dressed annihilators
\begin{equation}
 Q_R=cD_R+s\widehat X_L,
 \qquad Q_L=cD_L+s\widehat X_R.
 \label{eq:ch19-finite-Q-def}
\end{equation}
Their matrices, with columns representing the action on basis vectors,
are
\begin{equation}
 Q_R=\begin{pmatrix}
 0&c&0\\ 0&0&0\\ s&0&0
 \end{pmatrix},\qquad
 Q_L=\begin{pmatrix}
 0&0&c\\ s&0&0\\ 0&0&0
 \end{pmatrix}.
 \label{eq:ch19-finite-Q-matrices}
\end{equation}
For real $c,s$, direct multiplication gives
\begin{equation}
 H_Q:=Q_R^\dagger Q_R+Q_L^\dagger Q_L
   =\operatorname{diag}(2s^2,c^2,c^2).
 \label{eq:ch19-compressed-gram}
\end{equation}
Choose the admissible hyperbolic parameters $c=5/3$, $s=4/3$.
Then $c^2-s^2=1$ but
\begin{equation}
 H_Q=\operatorname{diag}\!\left(\frac{32}{9},
                                    \frac{25}{9},\frac{25}{9}\right)
 \succeq\frac{25}{9}I.
 \label{eq:ch19-compressed-witness}
\end{equation}
For every $v\in\mathcal H_{\leq1}$,
\begin{equation}
 \|Q_Rv\|^2+\|Q_Lv\|^2
 =\langle v,H_Qv\rangle
 \geq\frac{25}{9}\|v\|^2.
 \label{eq:ch19-compressed-coercivity}
\end{equation}
Therefore $\ker Q_R\cap\ker Q_L=\{0\}$.  There is no normalized
vector $\psi$ satisfying both $Q_R\psi=Q_L\psi=0$.

The density-matrix formulation gives the same obstruction without
assuming that the state is pure.  If $\Gamma\geq0$ has trace one, then
\begin{equation}
 \operatorname{Tr}(\Gamma H_Q)
 \geq\frac{25}{9}\operatorname{Tr}\Gamma
 =\frac{25}{9}>0.
 \label{eq:ch19-no-compressed-vacuum}
\end{equation}
In fact
$\operatorname{Tr}(\Gamma Q_\nu^\dagger Q_\nu)
 =\|Q_\nu\Gamma^{1/2}\|_{\mathrm{HS}}^2$;
vanishing for both $\nu=R,L$ would force the support of $\Gamma$
into their common kernel, hence $\Gamma=0$, a contradiction.  The
accepted A08 result \texttt{compressed\_no\_vacuum\_density}
proves the stronger algebraic impossibility of a trace-one matrix
satisfying $Q_R\Gamma=Q_L\Gamma=0$, without even requiring positivity.

Yet $H_Q$ is a perfectly valid finite self-adjoint Hamiltonian.  Its
ground energy is $25/9$, with ground eigenspace
$\operatorname{span}\{|R\rangle,|L\rangle\}$.  The density matrix
\begin{equation}
 \Gamma_{\mathrm{gs}}
  =\frac12\bigl(|R\rangle\langle R|
                   +|L\rangle\langle L|\bigr)
  =\operatorname{diag}\!\left(0,\frac12,\frac12\right)
 \label{eq:ch19-finite-ground-density}
\end{equation}
is positive and normalized, and satisfies
$\langle H_Q\rangle_{\Gamma_{\mathrm{gs}},\mathrm{fin}}=25/9$.
It is a genuine ground state but is \emph{not} a joint annihilation
vacuum.  By contrast, when $s=0$ the vector $|0\rangle$ belongs to
$\ker Q_R\cap\ker Q_L$.  Thus the obstruction in this example is
specific to the nonzero-mixing case; it does not say that every finite
Hamiltonian lacks a ground vector, nor that every compressed annihilator
family has trivial common kernel.

The same conclusion arises in a less artificial setting from the
finite-Fock analysis of the accepted Chapter~19 implementation.
There, under the explicitly required positivity of the band parameter,
admissibility of the chosen charge sector, positive mode indices inside
the cutoff, and both branch-margin hypotheses, the finite undressed
ground ket need not be the interacting dressed vacuum when $s\ne0$.
The theorem \texttt{ground\_not\_dressed\_vacuum} supplies a concrete
witness.  It does not construct a different finite interacting ground
state; that would require a specified finite Hamiltonian and its own
spectral analysis.  The three-dimensional model above is a diagnostic
counterexample, not a substitute for that analysis.

\subsection{What is established, and what remains to be compared}

There are therefore two well-defined, but mathematically different,
evaluation procedures:
\begin{equation}
 \underbrace{\langle X\rangle_{\psi,\mathrm{fin}}
 \quad\text{or}\quad
 \langle X\rangle_{\Gamma,\mathrm{fin}}}_{
 \substack{\text{finite CAR or finite}\text{ compressed Hamiltonian model}}}
 \qquad\text{versus}\qquad
 \underbrace{\omega_P(a)}_{
 \substack{\text{uncompressed abstract}\text{ CCR polynomial model}}}.
 \label{eq:ch19-two-state-layers}
\end{equation}
The first evaluates operators on a specified finite Hilbert space.
The second evaluates elements of an uncompressed relation algebra and
is positive by its polynomial construction and inverse Bogoliubov
pullback.  Both are normalized positive \emph{linear} functionals;
their domains, operator relations, and vacuum conditions differ.
There is no canonical equality between them.

An actual comparison would require, at minimum, a specified map taking
chosen abstract observables into finite-CAR observables (or a common
comparison algebra), a chosen finite state, the precise mode and energy
budgets, and an estimate of the errors caused by projection and finite
band edges.  Pointwise identities on a controlled excitation sector do
not automatically extend to all operator products; any estimate must
keep track of the excursions of intermediate states.  Subsequent
sections first compute ordered contractions using $\omega_P$, then
construct their finite spatial kernels, and finally address which
finite-Fock comparisons the available hypotheses justify.  No continuum
power-law or physical finite-ground-state equality is claimed here.
